\section{Reproducibility and Operational Implications}
\label{sec:discussion}

\subsection{RQ9: Can the Decisions and Numbers Be Checked?}
The compact evidence package contains per-target certification failures, both confidence allocations, selected shifts, paired evaluation arrays, cost components, the fresh-query summary, and generated table inputs. One script regenerates figures and typeset rows; another reanalyzes frozen native-response files through explicit repository and replay-root arguments. It checks the original twenty target decisions before deriving the new sensitivities. Historical records and the prospective source-slack stage remain separately labeled.

The preceding artifact audit preserved 80 pre-audit scientific files, passed a neutral-path smoke and 28 focused tests, and exposed guarded replay entry points for DARTH, Vamana, and Deep1M. In a clean Python environment built from the anonymous allowlist alone, the manuscript package passes 244 earlier evidence checks, reproduces all generated table text exactly, regenerates six figures, and verifies 210/210 source-only fresh-query checks. The post-seal checker additionally validates all 1,104 target decisions, 48 target summaries, four cost summaries, and both cost ledgers. It includes sanitized registration records for both fresh-query stages. These are reference-environment checks, not independent native-search replication. Recreating native searches still requires separately provisioned vectors, truth, indexes, binaries, and role manifests. The supplement gives the statistical recipes and detailed certificate tables; the package maps each claim to its compact evidence source.

\paragraph{Interpretation for index maintenance.}
The experiments suggest three operator decisions. After a graph-only rebuild, a historical budget is evidence about a source response, not the new one. After refresh, test the unchanged policy on the target rather than interpreting high mean recall as query-risk qualification. If recovery is considered, evaluate its completed candidate/fallback procedure and the horizon needed to pay for information. ICBA makes these checks explicit; TCP demonstrates target-calibrated recovery on a recurring workload, while the fresh-query bridge separates low-cost source candidate construction from target qualification.

\paragraph{Remaining boundaries.}
The graph-only result is a retrospective response diagnostic; TCP's positive evidence is a Faiss-HNSW recurring-profile case on two 100K datasets. The source-only boundary is prospective only for its original 500/500 roles. The post-seal replication uses new 500/500 target roles, remains conditional on 24 registered Faiss builds, and does not establish cross-implementation recovery. Its target-stage NDC ledger is complete, but its historical source-acquisition and full wall-clock ledgers are not estimable. Deep1M extends response evidence, not recovery economics. The TCP certification sensitivity is post-hoc and per target. The experiments do not identify a causal navigation-edge mechanism, demonstrate a learned unseen-query predictor, or establish hardware-general latency gains. Dynamic graph repair and target-trained adaptive methods are alternatives, not defeated baselines.

\section{Conclusion}
Search budgets are part of an index's operational state. On unchanged collections, source-derived first-passing actions incur substantial additional failure across hnswlib and Faiss rebuilds. After refresh, an unchanged recurring-query policy deteriorates and TCP recovers useful work on a recurring-query population. On fresh queries, a simpler source-derived slack rule retains work only on the registered Faiss grids; a post-seal replication then shows that the fixed Faiss candidate can be target-certified without losing its mean or tail benefit. ICBA makes these outcomes comparable without conflating them: certification, serving-work savings, target-stage amortization, and complete lifecycle value remain distinct. Each must be decided for the target index, implementation, action grid, and workload.
